% --- Archon physics formalization source begin ---
% archon:physics
% archon:covers PhyXMiniProblems/problem_phyx_mini_0106.lean
% archon:source-report reports/phyx_mini/problem_phyx_mini_0106.source.json

\chapter{Physics problem phyx\_mini\_0106}
\label{ch:PhyXMiniProblems_problem_phyx_mini_0106}

\paragraph{Problem source.}
\#\# Physical scenario

An object surrounded by a translucent material with index of refraction greater than unity appears enlarged. This includes items such as plants or shells inside plastic or resin paperweights, as well as ancient bugs trapped in amber. Consider a spherical object in air with radius \$r\$ fixed at the center of a sphere with radius \$R > r\$ with refractive index \$n\$, as shown in figure. When viewed from outside, the spherical object appears to have radius \$r' > r\$.The object is a spherical dandelion seed head with diameter \$45.0 \textbackslash{}, 	ext\{mm\}\$ fixed at the center of a solidified resin sphere with radius \$80.0 \textbackslash{}, 	ext\{mm\}\$ and index of refraction \$1.53\$.

\#\# Question

What diameter will the seed head appear to have when viewed from outside?

\#\# Answer choices

- A: A: 68.9mm

- B: B: 67.6mm

- C: C: 66.7mm

- D: D: 69.2mm

\#\# Figure caption (auxiliary; use the image as primary evidence)

The image depicts a diagram with two concentric circles and various labeled elements:

1. **Circles:**
   - A smaller orange circle labeled with radius \textbackslash{}( r' \textbackslash{}).
   - A larger blue circle labeled with an index \textbackslash{}( n \textbackslash{}).

2. **Points:**
   - A point on the circumference of the orange circle labeled \textbackslash{}( r' \textbackslash{}).
   - A point on the circumference of the blue circle connecting to the orange circle.

3. **Arrows and Lines:**
   - Two purple arrows extend outward from the circles.
   - A purple line connects the center of the orange circle to a point on the larger circle.
   - Dashed lines indicate angles and radii, including:
     - \textbackslash{}( \textbackslash{}theta\_a \textbackslash{}), an angle between the radius \textbackslash{}( R \textbackslash{}) and the outgoing arrow.
     - \textbackslash{}( \textbackslash{}theta\_b \textbackslash{}), an angle related to the positioned arrows.

4. **Text and Labels:**
   - \textbackslash{}( R \textbackslash{}): Represents the radius from the center of the orange circle to a point on the blue circle.
   - \textbackslash{}( r \textbackslash{}): A marked radius line inside the smaller circle (orange).
   - \textbackslash{}( \textbackslash{}theta\_a \textbackslash{}) and \textbackslash{}( \textbackslash{}theta\_b \textbackslash{}): Angles associated with the radii and arrows.

The diagram seems to illustrate geometric or trigonometric relations, possibly related to optics or wave propagation.

\#\# Dataset metadata

Geometrical Optics; Physical Model Grounding Reasoning, Spatial Relation Reasoning

\paragraph{Recorded answer/context.}
A: A: 68.9mm

\paragraph{Figure/image path.}
/root/proposal\_for\_physic/science-mango/phyx\_mini\_run/phyx\_data/test\_image/106.png

\paragraph{Formalization target.}
create a compiling Lean file with sorry bodies at `PhyXMiniProblems/problem\_phyx\_mini\_0106.lean`. The Lean declarations must preserve the physical quantities, dimensions or dimensional roles, figure labels, governing-law hypotheses, and final relation expressed by this problem.
Use Mathlib/Physlib names found through LeanExplore where available. If a physics API is missing, introduce faithful local abstractions rather than scalar placeholder aliases.

\begin{theorem}[Physics formalization target]
\label{thm:physics:phyx_mini_0106:target}
\lean{PhyXMiniProblems.ProblemPhyXMini0106.seed_head_apparent_diameter_matches_answer_A}
\uses{def:physics:phyx-mini-0106:phyxminiproblems-problemphyxmini0106-millimetersvalue, def:physics:phyx-mini-0106:phyxminiproblems-problemphyxmini0106-sphericalresinviewingsetup, def:physics:phyx-mini-0106:phyxminiproblems-problemphyxmini0106-hasproblemreadouts, def:physics:phyx-mini-0106:phyxminiproblems-problemphyxmini0106-matchesconcentricspherefigure, def:physics:phyx-mini-0106:phyxminiproblems-problemphyxmini0106-satisfiessphericalrefractionlaws, def:physics:phyx-mini-0106:phyxminiproblems-problemphyxmini0106-matchesanswertonearesttenth}
The assigned autoformalize agent should translate this physics problem into Lean declarations in the covered file, with theorem and lemma proof bodies written as `by sorry`.
\end{theorem}
\begin{proof}
This is an autoformalization task, not a proof task. Produce statements that can later be proved by the physics prover without weakening the physical model.
\end{proof}

% --- Archon declaration topology begin ---
\section{Declaration topology}
\begin{definition}[Length Quantity]
  \label{def:physics:phyx-mini-0106:phyxminiproblems-problemphyxmini0106-lengthquantity}
  \lean{PhyXMiniProblems.ProblemPhyXMini0106.LengthQuantity}
  A physical length represented coherently in every choice of units
\end{definition}
\begin{proof}
  This is the declared model definition of Length Quantity.
\end{proof}
\begin{definition}[millimeter Unit Choices]
  \label{def:physics:phyx-mini-0106:phyxminiproblems-problemphyxmini0106-millimeterunitchoices}
  \lean{PhyXMiniProblems.ProblemPhyXMini0106.millimeterUnitChoices}
  SI unit choices with length readouts expressed in millimetres
\end{definition}
\begin{proof}
  This is the declared model definition of millimeter Unit Choices.
\end{proof}
\begin{definition}[millimeters Value]
  \label{def:physics:phyx-mini-0106:phyxminiproblems-problemphyxmini0106-millimetersvalue}
  \lean{PhyXMiniProblems.ProblemPhyXMini0106.millimetersValue}
  \uses{def:physics:phyx-mini-0106:phyxminiproblems-problemphyxmini0106-lengthquantity, def:physics:phyx-mini-0106:phyxminiproblems-problemphyxmini0106-millimeterunitchoices}
  Read a physical length as a real number of millimetres
\end{definition}
\begin{proof}
  This is the declared model definition of millimeters Value.
\end{proof}
\begin{definition}[Sphere Label]
  \label{def:physics:phyx-mini-0106:phyxminiproblems-problemphyxmini0106-spherelabel}
  \lean{PhyXMiniProblems.ProblemPhyXMini0106.SphereLabel}
  The two physical spheres whose common centre is marked in the figure
\end{definition}
\begin{proof}
  This is the declared model definition of Sphere Label.
\end{proof}
\begin{definition}[Optical Medium]
  \label{def:physics:phyx-mini-0106:phyxminiproblems-problemphyxmini0106-opticalmedium}
  \lean{PhyXMiniProblems.ProblemPhyXMini0106.OpticalMedium}
  The optical media on the two sides of the spherical outer interface
\end{definition}
\begin{proof}
  This is the declared model definition of Optical Medium.
\end{proof}
\begin{definition}[Figure Ray]
  \label{def:physics:phyx-mini-0106:phyxminiproblems-problemphyxmini0106-figureray}
  \lean{PhyXMiniProblems.ProblemPhyXMini0106.FigureRay}
  The symmetric tangent rays drawn above and below the common centre
\end{definition}
\begin{proof}
  This is the declared model definition of Figure Ray.
\end{proof}
\begin{definition}[Spherical Resin Viewing Setup]
  \label{def:physics:phyx-mini-0106:phyxminiproblems-problemphyxmini0106-sphericalresinviewingsetup}
  \lean{PhyXMiniProblems.ProblemPhyXMini0106.SphericalResinViewingSetup}
  \uses{def:physics:phyx-mini-0106:phyxminiproblems-problemphyxmini0106-lengthquantity, def:physics:phyx-mini-0106:phyxminiproblems-problemphyxmini0106-spherelabel, def:physics:phyx-mini-0106:phyxminiproblems-problemphyxmini0106-opticalmedium, def:physics:phyx-mini-0106:phyxminiproblems-problemphyxmini0106-figureray}
  Physical quantities and labelled incidence geometry for the spherical resin viewing setup. thetaA is the angle inside the resin between the upper ray and the spherical-interface normal; thetaB is the refracted angle in air. The abstract geometric relations retain the physical roles of concentricity, tangency, and the parallel emergent rays without replacing those notions by untyped scalar flags
\end{definition}
\begin{proof}
  This is the declared model definition of Spherical Resin Viewing Setup.
\end{proof}
\begin{definition}[Has Problem Readouts]
  \label{def:physics:phyx-mini-0106:phyxminiproblems-problemphyxmini0106-hasproblemreadouts}
  \lean{PhyXMiniProblems.ProblemPhyXMini0106.HasProblemReadouts}
  \uses{def:physics:phyx-mini-0106:phyxminiproblems-problemphyxmini0106-millimetersvalue, def:physics:phyx-mini-0106:phyxminiproblems-problemphyxmini0106-sphericalresinviewingsetup}
  Numerical readouts stated in the problem. The seed-head diameter and resin radius are millimetre readouts, whereas refractive indices are dimensionless. No apparent radius or apparent diameter occurs in these data
\end{definition}
\begin{proof}
  This is the declared model definition of Has Problem Readouts.
\end{proof}
\begin{definition}[Matches Concentric Sphere Figure]
  \label{def:physics:phyx-mini-0106:phyxminiproblems-problemphyxmini0106-matchesconcentricspherefigure}
  \lean{PhyXMiniProblems.ProblemPhyXMini0106.MatchesConcentricSphereFigure}
  \uses{def:physics:phyx-mini-0106:phyxminiproblems-problemphyxmini0106-sphericalresinviewingsetup}
  Qualitative and ray-geometric information read from the concentric-sphere figure. In particular, the object lies strictly inside the resin sphere and its apparent image is enlarged. The sine inequalities select the physical acute-angle branch of the displayed refraction geometry
\end{definition}
\begin{proof}
  This is the declared model definition of Matches Concentric Sphere Figure.
\end{proof}
\begin{definition}[Satisfies Spherical Refraction Laws]
  \label{def:physics:phyx-mini-0106:phyxminiproblems-problemphyxmini0106-satisfiessphericalrefractionlaws}
  \lean{PhyXMiniProblems.ProblemPhyXMini0106.SatisfiesSphericalRefractionLaws}
  \uses{def:physics:phyx-mini-0106:phyxminiproblems-problemphyxmini0106-sphericalresinviewingsetup}
  The governing geometry and refraction laws for the ray construction. For every unit choice, tangency of the internal ray gives R sin(thetaA) = r. Snell's law at the resin--air interface gives n\_resin sin(thetaA) = n\_air sin(thetaB). Since the emergent ray is parallel to the viewing axis, its height is the apparent radius r' = R sin(thetaB). The remaining two laws relate each diameter to its radius. None of these laws assigns a numerical value to the requested apparent diameter
\end{definition}
\begin{proof}
  This is the declared model definition of Satisfies Spherical Refraction Laws.
\end{proof}
\begin{definition}[Answer Choice]
  \label{def:physics:phyx-mini-0106:phyxminiproblems-problemphyxmini0106-answerchoice}
  \lean{PhyXMiniProblems.ProblemPhyXMini0106.AnswerChoice}
  Labels of the four displayed multiple-choice answers
\end{definition}
\begin{proof}
  This is the declared model definition of Answer Choice.
\end{proof}
\begin{definition}[diameter Millimeters]
  \label{def:physics:phyx-mini-0106:phyxminiproblems-problemphyxmini0106-answerchoice-diametermillimeters}
  \lean{PhyXMiniProblems.ProblemPhyXMini0106.AnswerChoice.diameterMillimeters}
  \uses{def:physics:phyx-mini-0106:phyxminiproblems-problemphyxmini0106-answerchoice}
  Apparent-diameter readout in millimetres printed for each answer choice
\end{definition}
\begin{proof}
  This is the declared model definition of diameter Millimeters.
\end{proof}
\begin{definition}[Matches Answer To Nearest Tenth]
  \label{def:physics:phyx-mini-0106:phyxminiproblems-problemphyxmini0106-matchesanswertonearesttenth}
  \lean{PhyXMiniProblems.ProblemPhyXMini0106.MatchesAnswerToNearestTenth}
  \uses{def:physics:phyx-mini-0106:phyxminiproblems-problemphyxmini0106-lengthquantity, def:physics:phyx-mini-0106:phyxminiproblems-problemphyxmini0106-millimetersvalue, def:physics:phyx-mini-0106:phyxminiproblems-problemphyxmini0106-answerchoice}
  A physical diameter matches a choice displayed to the nearest tenth of a millimetre when its readout differs by at most 0.05 mm
\end{definition}
\begin{proof}
  This is the declared model definition of Matches Answer To Nearest Tenth.
\end{proof}
\begin{definition}[recorded Dataset Answer]
  \label{def:physics:phyx-mini-0106:phyxminiproblems-problemphyxmini0106-recordeddatasetanswer}
  \lean{PhyXMiniProblems.ProblemPhyXMini0106.recordedDatasetAnswer}
  \uses{def:physics:phyx-mini-0106:phyxminiproblems-problemphyxmini0106-answerchoice}
  The answer label recorded in the supplied dataset
\end{definition}
\begin{proof}
  This is the declared model definition of recorded Dataset Answer.
\end{proof}
% --- Archon declaration topology end ---
% --- Archon physics formalization source end ---
